\begin{question}[index=17,points=10]
求函数 $f(x,y)=(x-y)\eu^{-x^2-y^2}$ 的极值.
\end{question}
\QuestionBreak\WorkSpace

\begin{question}[index=18,points=12]
设 $y=f(x)=\dfrac12x^2$ 与 $y=g(x)=\sqrt{x-\dfrac34}$ 相切于点 $P$，$y=g(x)$ 与 $x$ 轴交于点 $A$，$O$ 为坐标原点.
\par（Ⅰ）求曲边图形 $OPA$ 绕 $y$ 轴旋转一周所得旋转体的体积；
\par（Ⅱ）求曲边图形 $OPA$ 绕 $x$ 轴旋转一周所得旋转体的侧面积.
\end{question}
\QuestionBreak\WorkSpace\WorkPageBreak

\begin{question}[index=19,points=12]
设函数 $z=f(x,y)$ 由方程 $F(x+z,y+z)=0$ 确定，其中 $F$ 为可微函数，计算
\[I=\iint_D\frac{xy}{y^2-x^2-1}\left(\frac{\partial z}{\partial x}+\frac{\partial z}{\partial y}\right)\dd x\dd y,\quad D=\{(x,y)\mid x^2+y^2\le1,0\le y\le x\}.\]
\end{question}
\QuestionBreak\WorkSpace

\begin{question}[index=20,points=12]
设二阶线性微分方程 $(1-x^2)\dfrac{\dd^2y}{\dd x^2}+(2\sqrt{1-x^2}-x)\dfrac{\dd y}{\dd x}+y=2x$.
\par（Ⅰ）作变换 $x=\sin t$（$-\dfrac\pi2\le t\le\dfrac\pi2$），把原方程变换成 $y$ 关于 $t$ 的微分方程；
\par（Ⅱ）求原方程的通解.
\end{question}
\QuestionBreak\WorkSpace\WorkPageBreak

\begin{question}[index=21,points=12]
设正定二次型 $f(x_1,x_2,x_3)=2x_1^2+ax_2^2+ax_3^2-4x_2x_3$ 在正交变换 $\vect X=\vect Q\vect Y$ 下的标准形为 $\lambda(y_1^2+y_2^2)+\mu y_3^2$.
\par（Ⅰ）求正交矩阵 $\vect Q$；
\par（Ⅱ）若存在可逆矩阵 $\vect P$，对任意 $\vect X=(x_1,x_2,x_3)^{\mathrm T}$，有 $f(x_1,x_2,x_3)=|\vect P\vect X|^2$，求 $\vect P$.
\end{question}
\QuestionBreak\WorkSpace

\begin{question}[index=22,points=12]
设随机变量 $X_1$ 与 $X_2$ 相互独立，且均服从 $N(0,1)$．记 $X=\max(X_1,X_2)$，$Y=\min(X_1,X_2)$，$Z=X-Y$.
\par（Ⅰ）求 $Z$ 的概率密度 $f_Z(z)$ 和 $E(Z)$；
\par（Ⅱ）求二维随机变量 $(X,Y)$ 的分布函数.
\end{question}
\WorkSpace
